\section{Background}

\label{sec:background}

\subsection{Group-Relative Policy Optimization}

GRPO~\cite{shao2024deepseekmath} extends the REINFORCE
algorithm~\cite{williams1992simple} by computing advantages relative to a
group of simultaneously generated rollouts, eliminating the need for a learned
value baseline. Given task $x$, a group of $G$ rollouts
$\{y^{(i)}\}_{i=1}^G \sim \pi_\theta(\cdot \mid x)$, and associated scalar
rewards $\{r^{(i)}\}_{i=1}^G$, the group-relative advantage is:
\begin{equation}
    A^{(i)} = \frac{r^{(i)} - \mu_G}{\sigma_G + \epsilon},
    \quad \mu_G = \frac{1}{G}\sum_{j=1}^G r^{(j)},
    \quad \sigma_G = \sqrt{\frac{1}{G-1}\sum_{j=1}^G (r^{(j)} - \mu_G)^2},
    \label{eq:grpo-advantage}
\end{equation}
with $\epsilon = 10^{-8}$ for numerical stability. Standard GRPO then performs
a clipped policy gradient update on $\theta$:
\begin{equation}
    \mathcal{L}_{\text{GRPO}}(\theta) = -\frac{1}{G}\sum_{i=1}^G \min\!\left(
        \rho^{(i)} A^{(i)},\; \text{clip}(\rho^{(i)}, 1{-}\delta, 1{+}\delta) A^{(i)}
    \right) + \lambda_{\text{KL}}\,\KL\bigl[\pi_\theta \| \pi_{\text{ref}}\bigr],
    \label{eq:grpo-loss}
\end{equation}
where $\rho^{(i)} = \pi_\theta(y^{(i)} \mid x) / \pi_{\text{old}}(y^{(i)} \mid x)$
is the importance weight and $\pi_{\text{ref}}$ is a reference policy.

\subsection{Training-Free GRPO (TF-GRPO)}

TF-GRPO~\cite{hanzo2026aso} replaces the gradient update in
Equation~\eqref{eq:grpo-loss} with a semantic extraction and compression
pipeline. The key insight is that the advantages $\{A^{(i)}\}$ encode
\emph{what makes one solution better than another} in a form that can be
distilled into reusable natural-language patterns, called \emph{semantic
advantages}:
\begin{equation}
    S^{(i)} = \mathcal{F}_{\text{LLM}}\!\left(x,\;\{y^{(j)}, r^{(j)}\}_{j=1}^G,\; i\right),
    \label{eq:semantic-advantage}
\end{equation}
where $\mathcal{F}_{\text{LLM}}$ denotes an LLM-based extraction call that
produces a structured text description of why rollout $i$ outperformed or
underperformed the group mean.

These semantic advantages are then compressed into 1-bit token-level expert
factors and applied at decode time via a Product-of-Experts (PoE)
distribution~\cite{hinton2002training}:
\begin{equation}
    \pi_{\text{ASO}}(y_t \mid x, y_{<t}) \propto \pi_\theta(y_t \mid x, y_{<t})
    \prod_{m=1}^{M} \phi_m(y_t \mid x, y_{<t})^{\eta_m},
    \label{eq:poe}
\end{equation}
where $\eta_m = \log(q_m / (1 - q_m))$ and $q_m \in (0,1)$ is the quality
attestation score for prior $m$. The composite reward used in ASO is:
\begin{equation}
    g^{(i)} = \alpha\, r^{(i)} + \beta\, u^{(i)},
    \label{eq:aso-reward}
\end{equation}
where $r^{(i)}$ is the extrinsic reward (test pass rate) and $u^{(i)}$ is the
epistemic reward (solution diversity and novelty).

\subsection{The YouTou-Agent Three-Stage Pipeline}

\cite{tencent2025youtoagent} propose a three-stage pipeline for extracting
reusable experiences from groups of LLM trajectories:

\begin{description}[leftmargin=1.5em, font=\normalfont\bfseries]
    \item[Stage 1 -- Trajectory Summarization.]
    Each individual trajectory $\tau^{(i)}$ in a group is summarized into a
    structured description of the reasoning steps, decisions, and outcomes. The
    summary is constrained to at most 32 words and must identify the key
    decision that differentiated this trajectory from a baseline.

    \item[Stage 2 -- Group Advantage Extraction.]
    Given $G$ summaries from Stage~1, the group is analyzed as a unit. The
    extractor identifies patterns that distinguish high-reward trajectories from
    low-reward ones, and produces \emph{operations} on an experience library
    $\mathcal{E}$: \textsc{add} new experiences, \textsc{modify} existing ones,
    \textsc{delete} outdated entries, or \textsc{merge} near-duplicate entries.

    \item[Stage 3 -- Batch Consolidation.]
    Multiple Stage-2 outputs (from different task groups within a session) are
    consolidated to remove redundancies, resolve conflicts, and ensure no single
    experience exceeds 32 words. The final set of operations is applied atomically
    to $\mathcal{E}$.
\end{description}

The original YouTou-Agent framework applies this pipeline to text-only LLM
outputs. Agent-GRPO adapts all three stages to the richer structure of
tool-use trajectories.

\subsection{Agent Trajectories vs.\ Text Trajectories}
\label{sec:traj-comparison}

A text-only trajectory $\tau_{\text{text}} = (x, y, r)$ consists of a task
$x$, a text output $y$, and a scalar reward $r$. An agent trajectory carries
strictly more information:

\begin{itemize}[leftmargin=1.5em]
    \item \textbf{Tool selection sequence:} which tools were called, and in what
          order.
    \item \textbf{Parameter choices:} the arguments passed to each tool,
          reflecting the agent's understanding of the task.
    \item \textbf{Intermediate outputs:} the return values of each tool call,
          observable during execution.
    \item \textbf{Per-step success signals:} whether each tool call succeeded
          or raised an error, providing dense local reward information.
    \item \textbf{Recovery patterns:} how the agent responded to errors---by
          retrying, switching tools, or escalating.
\end{itemize}

This richer signal structure makes agent trajectories both more informative
and more complex to process. The three-stage pipeline must be extended to
handle sequences of (tool, parameters, output, success) tuples rather than
single text outputs.

% ============================================================================
% 3. AGENT-LEVEL GRPO
% ============================================================================
